% Two extended bodies, centre-of-mass position vectors, and a surface point.
% Reference: handdraw/two_body.png. Compile: xelatex two_body_tikz.tex
\documentclass[tikz,border=6pt]{standalone}
\usepackage{amsmath}
\usepackage[fontset=fandol]{ctex}
\usetikzlibrary{arrows.meta,calc}
\begin{document}
\begin{tikzpicture}[>=Stealth,thick,every node/.style={font=\fontsize{18}{22}\selectfont}]
  \coordinate (C1) at (-2,0);
  \coordinate (O) at (0,0);
  \coordinate (C2) at (4.4,0);
  \coordinate (P) at ($(C1)+(55:1.15)$);
  \draw[blue] (C1) circle (1.15);
  \draw[blue] (C2) circle (.58);
  \fill[blue] (O) circle (2.3pt);
  \fill[blue] (C2) circle (1.7pt);
  \draw[->] (O) -- (C1) node[midway,below] {$\vec{x}_1$};
  \draw[->] (O) -- (C2) node[midway,below] {$\vec{X}$};
  \node[below] at (C2) {$\vec{x}_2$};
  \draw[->] (C1) -- (P) node[midway,left] {$\vec{x}$};
  \draw[->] (C2) -- (P) node[midway,above,sloped] {$\vec{X}-\vec{x}$};
  \node[below] at (0,-.35) {$O$};
  \node[blue] at (-2,-1.9) {星体1};
  \node[blue] at (4.4,-1.9) {星体2};
\end{tikzpicture}
\end{document}
